\section{Method}
\label{sec:method}

% =========================================================================
% Core section. Target length for first draft: ~4 pages of drafted prose.
% This section is where the recursive paper substantively differs from
% the framework paper. The framework paper introduces the recursive
% schema fields as forward-compatibility hooks (\S 21.3) and gestures
% at the productive-versus-unproductive distinction in its discussion.
% Here we develop both in depth.
% =========================================================================

\subsection{Overview}
\label{sec:method-overview}

The companion framework paper's approximate $Q$-learning system fits a
parametric action-value function $\Qpsi$ over a structured-state MDP
whose schema endows every hypothesis, subgoal, and verifier call with a
stable identifier and with optional parent-pointer and children fields.
The collections themselves remain flat tuples, so the tree structure is
implicit in the identifier references rather than carried by nested
data. The framework paper exercises none of this in phase 1; it leaves
the recursive fields at their defaults and proceeds with a flat
deliberation surface. The present section specifies what changes when
those fields are exercised, and in particular how exercising them
reshapes the sample efficiency of the approximate $Q$-learning system
and the structure of the MCTS rollouts that produce its targets.

We argue four points. First, structural recursion has two qualitatively
different modes (Section~\ref{sec:method-productive}), one of which is
the engine of compositional reasoning and the other of which is
rumination, and the difference is operational: it depends on whether the
verifier signal can reach the leaves of the implicit tree and so
contribute usable Bellman targets to $\Qpsi$. Second, a bounded depth
limit is the structural prior on the MDP's transition function that
suppresses the rumination mode without foreclosing the compositional one
(Section~\ref{sec:method-depthlimit}). Third, the framework paper's
seven-verb action set $\Aset$ is sufficient for recursion at any depth,
because identifier targeting is independent of where the targeted entity
sits in the implicit tree, and the multi-head $Q$-function over the
fixed verb set stays simple regardless of state-space depth
(Section~\ref{sec:method-targeting}). Fourth, serializing the implicit
tree as flat tokens with depth markers lets a frozen LLM (which serves
as both encoder and stochastic transition model) resolve recursive
structure through native causal self-attention, with no graph neural
network and no tree-recursive cell
(Section~\ref{sec:method-serialization}). We close with how MCTS, as the
policy improvement operator over $\Qpsi$, interacts with recursive
state structure (Section~\ref{sec:method-mcts}) and a note on
canonicalization implications (Section~\ref{sec:method-canon}).

\subsection{Productive and Unproductive Recursion}
\label{sec:method-productive}

Structural recursion in our deliberation MDP confronts two failure modes
of the approximate $Q$-learning system. On one side, unbounded recursion
permits an agent to spend an arbitrary fraction of its decision budget
on critique of critique, where each nested level demands more
justification of the level above it without any independent verifier
signal beneath. Such trajectories are structurally legal under the
framework paper's schema: nothing in the type system prevents
\texttt{parent\_hypothesis} chains of arbitrary length, and the verb
\textsc{Decompose} is available whenever a hypothesis exists. The
terminal reward is zero or negative on such trajectories because no
\textsc{CommitAnswer} ever fires and the depth budget eventually
triggers the no-commit penalty, so a fully trained $\Qpsi$ will assign
them low value. The problem is sample efficiency: discovering that
meta-criticism is unproductive requires spending decision steps to
explore it, and those steps yield only the eventual no-commit signal as
a Bellman target. From the standpoint of the $Q$-learning update, an
unproductive trajectory contributes one weak terminal target spread
across many transitions, none of which carry informative intermediate
verifier signal.

On the other side, prohibiting recursion entirely (constraining
\texttt{parent\_hypothesis} to \texttt{None} always) loses expressivity.
A hypothesis like ``the integral over $[0, 1]$ equals $\tfrac{1}{2}$''
on a multi-step calculus problem may split naturally into ``the
antiderivative is $\tfrac{x^2}{2}$'' and ``substitution gives
$\tfrac{1}{2} - 0$,'' each of which is independently checkable. Without
the ability to attach the latter two as children of the former, the
agent must either flatten the decomposition (producing two unconnected
hypotheses that do not reflect their derivational relationship) or
commit the parent claim directly and forgo verifier checks on the
intermediates. Both options surrender something the recursive schema
was designed to support, and both deprive $\Qpsi$ of the informative
intermediate verifier signals that productive recursion makes
available.

The pattern that distinguishes the two failure modes is verifier reach.
We call a recursive trajectory \emph{productive} when terminal verifier
reward eventually enters at the leaves of the induced tree and
propagates through the MCTS backups into chance-node estimates
$\Qmcts$, which then serve as targets for $\Qpsi$. We call it
\emph{unproductive} when no verifier signal can be obtained at the
leaves within the depth budget. The framework paper's schema produces
both with equal facility; the trained $Q$-head must come to prefer the
productive mode, and the simulation budget must be large enough to
make that preference learnable from MCTS-improved targets.

\begin{definition}[Productive recursion]
\label{def:productive}
A recursive subtree rooted at hypothesis $h$ is \emph{syntactically
productive} (under a verifier $V$ and a depth budget $\depthlim$) if
there exists a finite trajectory from $h$ to a leaf, of depth at most
$\depthlim$, that terminates in a successful \textsc{InvokeVerifier}
call on a descendant hypothesis or in a \textsc{CommitAnswer} that
$V$ accepts. It is \emph{operationally productive} under a policy
$\pi$ if a syntactically productive trajectory is realized with
positive probability by the trajectory distribution that $\pi$
induces in the LLM-driven context MDP. The subtree is
\emph{unproductive} (in either sense) when the corresponding
quantifier fails.
\end{definition}

Both senses are needed in what follows. Syntactic productivity is a
property of the schema, the verifier, and the depth budget: it asks
whether \emph{any} reachable trajectory carries leaf-level verifier
signal. Operational productivity is the property the training loop
actually cares about: it asks whether the current policy $\pi$
\emph{realizes} such trajectories often enough to supply Bellman
targets. The two coincide in the limit (an asymptotically optimal
policy under verifier-anchored reward concentrates on syntactically
productive subtrees) but diverge during learning: a subtree may be
syntactically productive yet operationally inert because the early
policy never explores it, and a subtree may be syntactically
unproductive yet temporarily attractive because the policy's value
estimates have not yet converged. The depth limit
(Section~\ref{sec:method-depthlimit}) targets the syntactic
quantifier (it removes subtrees whose only paths to verifier signal
exceed $\depthlim$); the rumination hypothesis in the experimental
plan (Section~\ref{sec:experimental-plan}) targets the operational
quantifier (it asks whether the trained policy under R$\infty$ wastes
its simulation budget on operationally unproductive depth).

Productivity is defined relative to a verifier and a depth budget; it is
not a property of the recursion structure alone. The same syntactic tree
may be productive under one verifier (a symbolic algebra system that
accepts a wide class of equivalence-up-to-simplification claims) and
unproductive under another (a strict equality checker). The framework
paper's verifier-anchored reward principle (verifier signals are the
only source of Bellman targets) means productivity directly governs
reward, and so governs which subtrees $\Qpsi$ comes to favor as the
approximate $Q$-learning loop closes.

\paragraph{Implication for design.}
The architectural job, then, is not to ban one mode and admit the other.
The schema cannot tell them apart from syntax. The architectural job is
to make productive recursion cheap to explore and unproductive recursion
expensive to dwell on, so that the approximate $Q$-learning system
spends its samples on trajectories that yield informative Bellman
targets. The depth limit (Section~\ref{sec:method-depthlimit}) is the
principal mechanism, but it is not the only one: identifier-based action
targeting (Section~\ref{sec:method-targeting}) ensures that the verifier
can be invoked on any node in the tree regardless of depth, and the
serialization template (Section~\ref{sec:method-serialization}) ensures
that the LLM-as-encoder sees the implicit tree as legible structure
rather than as sequence noise.

\subsection{The Depth Limit as Structural Prior}
\label{sec:method-depthlimit}

The framework paper introduces a per-hypothesis depth limit
$\depthlim$ (called \texttt{tree\_depth\_limit} in the schema and set to
$3$ or $4$ by default) and uses it to bound the recursion exercise in
phase 2. The brevity of that introduction understates how load-bearing
the limit is. We develop here a fuller treatment, framing the depth
limit as a structural prior on the MDP's transition function rather
than as a search-time heuristic.

\begin{definition}[Tree depth]
\label{def:tree-depth}
The \emph{tree depth} of an entity (hypothesis, subgoal, or verifier
call) is the length of the longest parent-pointer chain reachable from
that entity by following \texttt{parent\_hypothesis} (or
\texttt{parent\_subgoal}, or \texttt{follow\_up\_to}) repeatedly until
a $\textsf{None}$ field is reached. The tree depth of a root entity
(one whose parent field is $\textsf{None}$) is zero. The tree depth of
a state $s$ is the maximum tree depth over all entities in $s$.
\end{definition}

The depth limit $\depthlim$ enters as a constraint on the
\textsc{Decompose} and \textsc{FormHypothesis} availability predicates
in the MDP itself: both are removed from the available action set
$\Aset(s)$ whenever the new entity they would create would exceed the
depth limit. The constraint is part of the transition function's
domain definition, not a soft penalty added to reward. The framework
paper's availability table is amended:

\begin{center}
\small
\begin{tabular}{@{}ll@{}}
\toprule
Verb & Additional depth condition \\
\midrule
\textsc{Decompose}       & target's tree depth $< \depthlim - 1$ \\
\textsc{FormHypothesis}  & if attached to parent, parent's depth $< \depthlim - 1$ \\
\textsc{CheckCase}       & target's tree depth $< \depthlim - 1$ \\
\textsc{ComputeSubexpr}  & no additional condition \\
\textsc{InvokeVerifier}  & no additional condition \\
\textsc{Backtrack}       & no additional condition \\
\textsc{CommitAnswer}    & no additional condition \\
\bottomrule
\end{tabular}
\end{center}

The three structurally recursive verbs are masked at the limit; the
remaining four are not, so the agent can always invoke the verifier on a
leaf hypothesis, abandon a subtree, or commit an answer regardless of
how deep the current state's tree is. The choice is deliberate: the
depth limit constrains how much structure the agent can build, not how
freely it can resolve or back out of existing structure. From the
$Q$-learning perspective, the limit shapes the support of $\Aset(s)$
and so shapes which transitions $\Qpsi$ ever sees during training.

\paragraph{Why a hard mask.}
A soft-penalty alternative (assign negative reward proportional to depth
above $\depthlim$) was considered and rejected for two reasons. First,
soft penalties interact with the verifier-anchored reward principle in
ways we want to avoid: the framework paper reserves Bellman targets for
verifier signals, and a depth penalty is exactly the kind of
aesthetically motivated additive reward that the framework paper warns
against. Mixing aesthetic and verifier rewards would muddy the Bellman
targets that train $\Qpsi$. Second, a soft penalty makes the policy
learnable but the state space exploration unbounded: an MCTS simulation
could still descend arbitrarily deep, paying the penalty at each step
but accruing no other information. The hard mask cuts off the descent
in the MDP itself and forces the agent to either resolve the existing
structure (via verifier, backtrack, or commit) or do something else.

\paragraph{Interaction with reward signal density.}
The depth limit interacts with the reward signal density in a way worth
making explicit. A trajectory of depth $d$ produces at most $d$
verifier-anchored events between root and leaf (one per
\textsc{InvokeVerifier} call on the path). Beyond that, additional
depth produces no new verifier signal; it can only re-derive,
re-decompose, or re-verify. Empirically (in the framework paper's
domain of math word problems and competition mathematics) the natural
decomposition depth for human-written solutions is rarely above three.
Setting $\depthlim = 3$ to $4$ matches this distribution and yields what
we expect to be the productive recursion regime, in which $\Qpsi$
receives a dense supply of verifier-anchored Bellman targets across the
implicit tree. Setting $\depthlim = 0$ disables recursion entirely and
provides the framework paper's phase 1 behavior. Setting
$\depthlim = \infty$ admits both modes and is the natural baseline for
measuring the depth-limit mechanism's contribution.

\paragraph{Adaptive depth limits.}
The depth limit can in principle vary by problem, by subtree, or by
policy state. The simplest version is global and fixed; this is what
the ablation in Section~\ref{sec:experimental-plan} tests. A
per-problem depth limit (learned from a difficulty estimator) and a
per-subtree depth limit (annealed as confidence in the subtree's
hypothesis grows) are natural extensions, but they add architectural
machinery beyond what the present paper needs to evaluate the core
claim. We defer them.

\subsection{Action Targeting at Arbitrary Depth}
\label{sec:method-targeting}

The framework paper's action set $\Aset$ is a fixed seven-element
categorical set: \textsc{Decompose}, \textsc{ComputeSubexpr},
\textsc{FormHypothesis}, \textsc{CheckCase}, \textsc{InvokeVerifier},
\textsc{Backtrack}, \textsc{CommitAnswer}. The multi-head $Q$-function
$\Qpsi$ assigns a scalar value to each verb at every state; MCTS, when
present, branches over these seven verbs at every decision node. The
target of a verb (which hypothesis to decompose, which fact to verify,
which subgoal to abandon) is supplied by the verb-instantiator as part
of its generated content, conditioned on the serialized state.

This separation between verb (categorical, the support of $\Qpsi$) and
target (continuous and id-valued, generated by the LLM) is the
mechanism by which the action space stays bounded as the state grows
tree-structured. We highlight three consequences.

\paragraph{No new verbs are required for recursion.}
\textsc{Decompose}\texttt{(hyp\_id=hyp\_4)} decomposes hypothesis
\texttt{hyp\_4}, regardless of whether \texttt{hyp\_4} is a top-level
hypothesis (\texttt{parent\_hypothesis = None}) or sits at depth 2 in
the implicit tree. The verb and its semantics are unchanged; the
verb-instantiator's targeting decision is the only thing that varies.
Similarly, \textsc{InvokeVerifier} can be invoked on
\texttt{hyp\_id=hyp\_7} where \texttt{hyp\_7} is a leaf at depth 3, with
exactly the same semantics as invoking it on a root hypothesis. The
$Q$-head's seven outputs stay the same regardless of state-space depth.

\paragraph{Branching factor at decision nodes is bounded.}
PUCT's chance-node statistics aggregate over a fixed seven-element
action space. The fact that the verb-instantiator generates a different
identifier in different LLM samples produces transition stochasticity,
which the framework paper handles via chance-node bifurcation, not
action-space inflation. The branching factor at decision nodes stays at
seven (or fewer, after availability masking) regardless of how deep the
tree gets. This preserves PUCT's reliance on a fixed action set without
needing progressive widening machinery, and it keeps $\Qpsi$'s
multi-head structure simple regardless of how rich the underlying state
becomes.

\paragraph{Verb-instantiator must produce valid ids.}
A subtle requirement: when the policy selects, say,
\textsc{Decompose}, the verb-instantiator must produce a continuation
that names an existing hypothesis or subgoal in the current state. The
serialization template (Section~\ref{sec:method-serialization}) makes
all existing identifiers visible in the LLM's prefix, and the
verb-specific prompt explicitly enumerates them. The free-form
continuation that follows is constrained by either constrained decoding
\citep{willard2023outlines} or a retry-with-correction loop. Failure
to produce a valid identifier returns control to the search with the
action masked as if it had not been chosen; this is the same machinery
the framework paper uses for \textsc{CommitAnswer} format violations.

\subsection{Encoding Recursive State without Recursive Encoders}
\label{sec:method-serialization}

The framework paper's serialization template represents the state as a
flat token sequence with explicit structural markers and depth-indexed
identifiers. We expand on this design choice here because recursion
puts it under more strain than the framework paper's first-version
flat-tree usage tests. The serialization is consumed by the frozen LLM
that serves as both encoder (its terminal hidden state feeds $\Qpsi$
and $\ptheta$) and stochastic transition model (its sampled
continuations realize the verb-instantiated next state).

\paragraph{The template.}
A state with implicit tree structure is serialized as follows
(simplified for exposition):

\begin{lstlisting}[basicstyle=\ttfamily\footnotesize]
[PROBLEM] {problem text} [/PROBLEM]
[FACT id=fact_1 source=given] {statement} [/FACT]
[OPEN_SUBGOAL id=sg_1 parent=none depth=0] {statement} [/OPEN_SUBGOAL]
[OPEN_SUBGOAL id=sg_2 parent=sg_1   depth=1] {statement} [/OPEN_SUBGOAL]
[HYPOTHESIS id=hyp_1 parent=sg_2     depth=1] {statement} [/HYPOTHESIS]
[HYPOTHESIS id=hyp_2 parent=hyp_1    depth=2] {statement} [/HYPOTHESIS]
[HYPOTHESIS id=hyp_3 parent=hyp_1    depth=2] {statement} [/HYPOTHESIS]
[VERIFIER_CALL id=vcall_1 target=hyp_3 verb=factor success=true]
   in:  {input}
   out: {output}
[/VERIFIER_CALL]
[COMMITTED_ANSWER] <none> [/COMMITTED_ANSWER]
[DEPTH] {turn count} [/DEPTH]
[STATE_END]
\end{lstlisting}

Three design choices distinguish this template from the framework
paper's flat-tree usage. First, the \texttt{parent=} field appears on
every entity capable of recursive nesting, with the literal token
\texttt{parent=none} for roots and the parent's identifier otherwise.
Second, the \texttt{depth=} field is rendered as an explicit integer
even though it is derivable from the parent chain. Third, sibling
hypotheses (\texttt{hyp\_2} and \texttt{hyp\_3} above) appear adjacent
in the serialization, ordered by trajectory time of introduction; the
parent-pointer is the load-bearing relational signal, not their
sequential position.

\paragraph{Why this suffices for a frozen LLM.}
The argument that causal self-attention can resolve this implicit tree
proceeds in three steps. First, the structural markers
(\texttt{[HYPOTHESIS id=hyp\_2 parent=hyp\_1 depth=2]} and so on) give
the LLM stable positional anchors for each entity. Decoder-only LLMs are
trained to make late-position hidden states integrate the entire
prefix, so the \texttt{[STATE\_END]} sentinel's hidden state has access
to every preceding token, including all parent-pointer references; this
is the hidden state that $\Qpsi$ reads off. Second, the LLM's attention
can match a hypothesis to its parent via identifier matching: when the
LLM is generating a continuation that references \texttt{hyp\_1},
attention heads can resolve which prior hypothesis introduced
\texttt{hyp\_1} by attending to the prior \texttt{id=hyp\_1} token. This
is standard in-context retrieval behavior of modern LLMs. Third, the
depth-indexed identifier (\texttt{depth=2}) provides redundant signal:
even if attention misses the parent-pointer chain, the depth integer
makes the recursion level explicit at the marker.

\paragraph{Comparison with explicit tree encoders.}
A tree-LSTM \citep{tai2015improved} would compose hidden states
recursively along the parent-child structure: each node's representation
is a function of its children's representations and its own content. A
graph attention network \citep{velickovic2018gat} would extend this to
arbitrary parent-pointer graphs via message passing. Both produce a
representation explicitly aware of the tree topology.

The frozen-LLM-plus-flat-serialization approach gives up that explicit
structural awareness. What it gains is architectural simplicity: no
extra encoder, no representation drift from a separate encoder family,
no fine-tuning required to make the encoder agree with the generator,
and a single hidden state on which $\Qpsi$ and $\ptheta$ both read. The
empirical question is whether modern LLMs' in-context attention is
a sufficient substitute. We have qualitative reasons to expect yes:
LLMs trained on code and structured text routinely resolve nested
braces, parent-child JSON references, and analogous structures via
attention alone, with no architectural tree-awareness. We have no
analytic guarantee. The ablation in
Section~\ref{sec:experimental-plan} addresses the question for our
specific deliberation MDP.

\paragraph{Pathological cases.}
Serialization breaks down in two regimes. First, if the implicit tree
becomes deep enough that distant parent-pointer references exit the
attention's effective receptive field (whether by raw distance or by
soft-attention dilution), the recursive structure becomes inaccessible.
Modern LLMs with $32$k or $128$k context windows make this a remote
concern for deliberation tasks where total context rarely exceeds a few
thousand tokens, but it is worth flagging. Second, if the implicit
graph becomes non-tree (a hypothesis referenced as parent by multiple
children, or cyclic via \texttt{follow\_up\_to}), the depth field loses
well-definedness and the serialization template needs amendment. Our
schema's type system rules out cycles structurally
(\texttt{parent\_hypothesis} is a non-recursive string identifier
field), and the framework paper's design avoids multi-parent structure
in phase 2. Phase-3 extensions to DAG structure are out of scope here.

\subsection{MCTS as Policy Improvement Over the Recursive State Space}
\label{sec:method-mcts}

In the framework paper, MCTS is the policy improvement operator over
the approximate $Q$-learning system: it uses $\Qpsi$ as leaf evaluator
and $\ptheta$ as PUCT prior, and produces $\Qmcts$ and $\pimcts$ as
targets back to the learned heads. Recursion changes what this
operator sees in two ways. The branching factor at decision nodes is
unaffected: PUCT still branches over the seven verbs in $\Aset(s)$.
The state space, however, becomes tree-structured: two states that
agree on flat-tuple contents but differ in the implicit tree of parent
pointers are distinct states, and MCTS treats them as such.

\paragraph{Tree reuse implications.}
The framework paper defers canonicalization to phase 2. With recursion
active, canonicalization becomes more delicate. Two states that reach
semantically equivalent trees (same hypotheses, same parent
relationships, possibly different identifier choices) ideally hash
identically for MCTS tree reuse and for $\Qpsi$ generalization. The
natural canonicalization is to rename identifiers along a canonical
traversal order (e.g., depth-first preorder), but this is more
sensitive to introduction order than the flat-tuple version: a
hypothesis introduced as \texttt{(parent=sg\_1, children=())}
canonicalizes differently from one introduced as
\texttt{(parent=sg\_1, children=(hyp\_5,))} even when the latter's
child is subsequently abandoned. We treat canonicalization as out of
scope for this paper and defer it to subsequent work; the ablation
grid below tests the depth-limit mechanism without relying on
canonicalization-based tree reuse.

\paragraph{Depth-versus-breadth tradeoff.}
The MCTS simulation budget $B$ at the root divides across the tree.
Deeper recursion concentrates the budget along long parent-pointer
chains; shallower recursion fans the budget across siblings. The
framework paper's $\depthlim$ default of $3$ to $4$ is a starting
point; the ablation in Section~\ref{sec:experimental-plan} measures
where on this tradeoff the productive regime sits empirically. Because
MCTS's chance-node statistics are precisely the source of training
targets for $\Qpsi$, this tradeoff also governs how informative each
unit of simulation budget is for the underlying approximate
$Q$-learning loop.

\paragraph{Backup along recursion structure.}
A subtle question: when MCTS backs up a verifier reward from a leaf
hypothesis \texttt{hyp\_3} (which is a child of \texttt{hyp\_1}), does
the reward flow through the chance-node parents (the
\textsc{InvokeVerifier} action that produced the verifier call) or
through the structural parents (the \texttt{parent\_hypothesis} link)?
In our design, only the former: MCTS's tree structure is the
decision-and-chance-node tree, not the implicit hypothesis tree.
Structural parents matter for the policy's representation of the state
(through serialization) and for the verifier's input (a verifier call
on \texttt{hyp\_3} may pull in \texttt{hyp\_1}'s statement as context),
but they are not channels for reward propagation in MCTS itself. This
keeps the MCTS algorithm unchanged from the framework paper; the
recursion lives entirely in the state representation that $\Qpsi$
and $\ptheta$ both ingest.

\subsection{Canonicalization and Idempotence}
\label{sec:method-canon}

% TODO: A brief subsection on the relationship between recursion and
% the canonicalization deferred from the framework paper. The main
% point: recursion makes the canonicalization problem harder (two
% structurally equivalent trees may differ in identifier choice and in
% introduction order of children), so the framework paper's deferral
% of canonicalization is, if anything, more justified once recursion is
% active. The implication for the approximate $Q$-learning system is
% that two semantically equivalent states may produce distinct Bellman
% targets, weakening $\Qpsi$'s generalization across structurally
% equivalent trees. We do not propose a canonicalization algorithm
% here; we flag the dependence between the two design choices.
